\documentclass[11pt]{article}
\usepackage[margin=1in]{geometry}
\usepackage{amsmath,amssymb}
\usepackage{graphicx}
\usepackage{booktabs}
\usepackage{xcolor}
\usepackage{caption}
\usepackage{float}
\usepackage[hidelinks]{hyperref}
\usepackage{enumitem}
\graphicspath{{figures/}{../figures/}}
% Robust include: renders figures/ or ../figures/ if present, else a labeled placeholder,
% so this compiles on Overleaf even before you upload the figures/ folder.
\newcommand{\figimg}[2]{%
  \IfFileExists{figures/#1.png}{\includegraphics[width=#2]{figures/#1}}{%
  \IfFileExists{../figures/#1.png}{\includegraphics[width=#2]{../figures/#1}}{%
  \fbox{\parbox[c][2.4cm][c]{#2}{\centering\ttfamily\footnotesize [#1.png]\\[2pt]\normalfont\itshape upload the repo's \texttt{figures/} folder to render}}}}}
\captionsetup{font=small}
\setlength{\parskip}{4pt}
\hypersetup{colorlinks=true,linkcolor=blue,urlcolor=blue,citecolor=blue}

\title{\textbf{facilitiesHelp.io}\\[2pt]
\large Evidence-Graded Trust, Medical-Desert Detection, and Causal Honesty\\ for 10{,}000 Messy Healthcare-Facility Records}
\author{Kumar Koushik Telaprolu\thanks{\texttt{kumarkoushik0805@gmail.com}} \quad\quad Chris Kwan\thanks{\texttt{ckwan48@gmail.com}}\\[2pt]
\small Databricks Apps \& Agents for Good 2026}
\date{}

\begin{document}
\maketitle

\begin{abstract}
\noindent A planner handed $10{,}088$ scraped Indian healthcare-facility records faces a trust problem: capability claims are unverified, key fields are sparse, and some values are impossible. We present \textbf{facilitiesHelp.io}, a Databricks App that treats every field as a \emph{claim to verify, not a fact}. We (i) grade every specialty a facility lists ($2{,}580$ distinct; $117{,}993$ graded claims) on evidence drawn from the facility's own text, attaching a calibrated $0\text{--}100$ confidence; (ii) validate that the grade reflects \emph{cited evidence rather than facility prominence} (a model predicting the grade from facility metadata alone achieves AUC $\approx 0.57$, near chance); (iii) classify district-level care gaps for $2{,}518$ capabilities, separating a \emph{real} gap from a \emph{data-poor} region; and (iv) build a causal layer on $706$ NFHS-5 districts that distinguishes real levers from coincidences. Our headline causal result: the sanitation--stunting association ($r=-0.51$) \emph{collapses to $\approx 0$ after adjusting for household wealth}, while female-schooling$\to$child-marriage ($r=-0.65$) and antenatal-care$\to$institutional-birth ($r=+0.60$) survive. Crucially, facility \emph{count} is only weakly linked to outcomes after adjustment, which reframes the supply--demand question for planners. Every claim cites its source text; uncertainty is communicated, never hidden.
\end{abstract}

\section{Introduction}
The challenge dataset contains $10{,}088$ Indian healthcare facilities described by $51$ columns: structured fields (name, location, controlled specialties) plus \emph{uneven free-text} descriptions of claimed capabilities, procedures, equipment, and services. Field coverage is highly non-uniform (Table~\ref{tab:coverage}), the free text is noisy and repetitive, and some structured values are impossible. The user---a non-technical planner, NGO coordinator, or analyst---must nonetheless make decisions they can trust: \emph{Can this facility do what it claims? Where are the real care gaps? Where should a patient go? What must be fixed before the data is trustworthy?}

Our contributions:
\begin{enumerate}[leftmargin=1.4em,itemsep=1pt]
\item A deterministic, auditable \textbf{trust-grading engine} over the facility's own evidence, with a calibrated confidence and a contradiction check (\S\ref{sec:trust}).
\item A \textbf{validation} that the grade is evidence-driven, not a popularity proxy (\S\ref{sec:valid}).
\item A \textbf{medical-desert} classifier that separates real gaps from data-poor regions for $2{,}518$ capabilities, addressing the supply--demand challenge (\S\ref{sec:desert}).
\item A \textbf{correlation--causation} layer on $706$ NFHS districts using structure learning, multi-estimator adjustment, and sensitivity analysis (\S\ref{sec:corr}--\ref{sec:causal}), with statistical-ML and geometric-DL extensions (\S\ref{sec:ml}).
\end{enumerate}

\section{Datasets and granular profiling}
We use three read-only Unity Catalog datasets:
\begin{itemize}[leftmargin=1.4em,itemsep=1pt]
\item \textbf{facilities} --- $10{,}088$ rows $\times\,51$ columns; the core, treated as claims to verify.
\item \textbf{india\_post\_pincode\_directory} --- $165{,}627$ post offices with district, state, latitude, longitude; powers geocoding and the facility$\to$district join.
\item \textbf{nfhs\_5\_district\_health\_indicators} --- $706$ districts $\times\,109$ indicators; powers the causal layer.
\end{itemize}

\paragraph{Coverage and missingness.} Substantive (non-empty content) coverage of the evidence fields is uneven (Table~\ref{tab:coverage}); $9{,}996/10{,}088$ records carry a postcode. We treat missingness honestly: sparse fields are shown as \emph{``not reported,''} never imputed, and the confidence of a claim reflects how much evidence actually exists.

\begin{table}[H]\centering\small
\begin{tabular}{lc|lc}
\toprule
Field & Coverage & Field & Coverage\\
\midrule
description & 100\% & numberDoctors & 36.4\%\\
capability & 99.7\% & capacity (beds) & 25.2\%\\
procedure & 92.5\% & yearEstablished & 47.8\%\\
equipment & 77.0\% & postcode & 99.0\%\\
\bottomrule
\end{tabular}
\caption{Evidence-field coverage. Useful signal is spread across \emph{description, capability, procedure, equipment, specialties}, and \emph{source\_urls}; all are claims, not ground truth.}
\label{tab:coverage}
\end{table}

\paragraph{Entity resolution.} The same facility is scraped from many sources. We deduplicate to one canonical record per pre-computed match key \texttt{cluster\_id} via
$\texttt{row\_number() OVER (PARTITION BY cluster\_id ORDER BY unique\_id)}$,
collapsing $10{,}088$ raw rows to $9{,}959$ facilities. A medallion pipeline (bronze$\to$silver$\to$gold) on a serverless SQL Warehouse then assembles the evidence text, counts independent source types, and joins facilities$\to$PIN$\to$district.

\section{The trust-grading engine}\label{sec:trust}
For each (facility, capability) we test two evidence signals against the facility's own text and count independent sources:
\[
\textsf{spec\_hit}=\mathbb{1}[\texttt{spec\_text} \sim r_{\text{spec}}],\quad
\textsf{text\_hit}=\mathbb{1}[\texttt{evidence\_text} \sim r_{\text{kw}}],\quad
n_{\text{src}}=|\text{distinct source\_types}|,
\]
where $\texttt{evidence\_text}=\text{lower}(\text{description}\,\Vert\,\text{capability}\,\Vert\,\text{procedure}\,\Vert\,\text{equipment})$ and $\sim$ denotes a regex match. The grade is
\[
\textsf{grade}=
\begin{cases}
\textsf{WEAK/SUSPICIOUS} & \text{facility type contradicts an acute claim}\\
\textsf{STRONG} & \textsf{spec\_hit}\wedge\textsf{text\_hit}\\
\textsf{PARTIAL} & \textsf{spec\_hit}\vee\textsf{text\_hit}\\
\textsf{NO CLAIM} & \text{otherwise (dropped)}
\end{cases}
\]
The contradiction rule flags a small \texttt{clinic/dentist/doctor/pharmacy} asserting an acute capability (ICU, oncology, $\dots$) in free text with \emph{no} matching specialty code. Confidence is a bounded, monotone function of corroboration:
\[
\textsf{conf}=
\begin{cases}
\min(95,\;70+5\,n_{\text{src}}) & \textsf{STRONG}\\
50 & \textsf{PARTIAL}\\
20 & \textsf{WEAK/SUSPICIOUS}
\end{cases}
\]
We apply this enhanced rubric (with the contradiction check) to $22$ high-stakes acute capabilities, and a generalized version to \emph{every} specialty a facility lists, yielding $117{,}993$ graded specialty-claims across $2{,}580$ specialties (Table~\ref{tab:grades}). Every grade is accompanied by the verbatim evidence and ranked source links, so a planner can audit it.

\begin{table}[H]\centering\small
\begin{tabular}{lcc}
\toprule
Grade & 22 acute capabilities & All $2{,}580$ specialties\\
\midrule
STRONG & $36{,}847$ & $40{,}151$\\
PARTIAL & $37{,}318$ & $64{,}898$\\
WEAK/SUSPICIOUS & $1{,}170$ & ---\\
CLAIMED (listed only) & --- & $12{,}944$\\
\midrule
Total & $75{,}335$ & $117{,}993$\\
\bottomrule
\end{tabular}
\caption{Graded-claim distribution. STRONG requires $\geq2$ independent evidence types and corroboration; most claims are PARTIAL, reflecting honest single-signal evidence.}
\label{tab:grades}
\end{table}

\begin{figure}[H]\centering
\figimg{capability_grades}{0.78\textwidth}
\caption{Evidence grades by capability. Maternity is best-evidenced; ICU/NICU are rarest---visibility, not necessarily absence.}
\label{fig:grades}
\end{figure}

\section{Validating the grade: reproducibility and non-popularity}\label{sec:valid}
Two questions: is the grade a faithful function of its evidence, and is it merely a proxy for facility prominence?
\begin{enumerate}[leftmargin=1.4em,itemsep=1pt]
\item \textbf{Reproducibility.} A logistic regression $\Pr(\textsf{STRONG})\sim \textsf{spec\_hit}+\textsf{text\_hit}+n_{\text{src}}$ attains $5$-fold ROC-AUC $=1.00$, confirming the grade is a transparent, deterministic function of its evidence (attribution $\approx 49\%$ specialty code, $\approx 51\%$ procedure/equipment text).
\item \textbf{Non-popularity.} A gradient-boosting classifier predicting \textsf{STRONG} from facility \emph{metadata only} (capacity, doctors, web presence, \#specialties, facility type)---\emph{excluding} the rule inputs---scores only $4$-fold ROC-AUC $=0.57$, barely above chance. Therefore the grade reflects \emph{cited evidence, not size or fame.}
\end{enumerate}
We additionally expose a SHAP-style per-claim attribution (Fig.~\ref{fig:fi}) that decomposes each confidence score into its contributing signals.

\begin{figure}[H]\centering
\figimg{ml_feature_importance}{0.7\textwidth}
\caption{Feature importance / attribution for the evidence model.}
\label{fig:fi}
\end{figure}

\section{Medical deserts and the supply--demand challenge}\label{sec:desert}
For each (capability, district) aggregated over distinct facilities we classify
\[
\textsf{class}=
\begin{cases}
\textsf{DATA-POOR} & |\{\text{facilities}\}|<3\\
\textsf{APPARENT CARE GAP} & n_{\textsf{STRONG}}<2\\
\textsf{EVIDENCED SUPPLY} & \text{otherwise.}
\end{cases}
\]
This is the honesty hinge: it separates a \emph{real} gap (facilities exist but little trustworthy evidence) from a region we simply \emph{under-measured} (Table~\ref{tab:gaps}). Pre-computed for all $2{,}518$ capabilities. That the modal cell is DATA-POOR is itself the finding---most district$\times$capability combinations cannot be judged from scraped data, and we say so rather than declaring ``no care.''

\begin{table}[H]\centering\small
\begin{tabular}{lc}
\toprule
District$\times$capability classification & Count\\
\midrule
DATA-POOR (can't tell) & $24{,}824$\\
EVIDENCED SUPPLY & $8{,}674$\\
APPARENT CARE GAP (real) & $4{,}675$\\
\bottomrule
\end{tabular}
\caption{Care-gap classification across $2{,}518$ capabilities and $706$ districts.}
\label{tab:gaps}
\end{table}

\paragraph{Supply vs.\ demand.} Facility data is \emph{supply} (what care exists, how trustworthy); NFHS is \emph{demand and outcomes}. They meet at the district key. The pivotal question---\emph{will adding facilities improve outcomes?}---is answered in \S\ref{sec:causal}: facility \emph{count} is only weakly linked to outcomes after adjustment, so a desert is necessary but not sufficient for an investment decision.

\section{Correlation analysis}\label{sec:corr}
We summarize linear association with the Pearson coefficient
\[
r=\frac{\sum_i (x_i-\bar x)(y_i-\bar y)}{\sqrt{\sum_i (x_i-\bar x)^2}\,\sqrt{\sum_i (y_i-\bar y)^2}}=\frac{\operatorname{cov}(X,Y)}{\sigma_X\,\sigma_Y}\in[-1,1],
\]
computed across the $706$ districts of \texttt{gold\_nfhs}; each row is a district and each column an NFHS percentage. Strength tiers: $|r|<0.1$ none, $0.1\text{--}0.3$ weak, $0.3\text{--}0.5$ moderate, $0.5\text{--}0.7$ strong, $>0.7$ very strong. Selected live values appear in Table~\ref{tab:corr}; the full matrix is the interactive heatmap (Fig.~\ref{fig:heat}).

\begin{table}[H]\centering\small
\begin{tabular}{lc}
\toprule
Pair (district level, $n=706$) & $r$\\
\midrule
sanitation $\leftrightarrow$ stunting & $-0.51$\\
clean fuel $\leftrightarrow$ stunting & $-0.42$\\
female schooling $\leftrightarrow$ child marriage & $-0.65$\\
ANC4 visits $\leftrightarrow$ institutional birth & $+0.60$\\
women overweight $\leftrightarrow$ high blood sugar & $+0.62$\\
\textbf{wealth} $\leftrightarrow$ sanitation & $+0.32$\\
\textbf{wealth} $\leftrightarrow$ stunting & $-0.43$\\
\bottomrule
\end{tabular}
\caption{Pearson correlations. The last two rows are the ``smoking gun'': wealth is positively linked to sanitation \emph{and} negatively to stunting, which mechanically manufactures the sanitation--stunting association.}
\label{tab:corr}
\end{table}

\begin{figure}[H]\centering
\figimg{corr_heatmap}{0.85\textwidth}
\caption{NFHS-5 correlation matrix across $706$ districts (red positive, blue negative).}
\label{fig:heat}
\end{figure}

\section{Causal inference: separating levers from coincidences}\label{sec:causal}
\paragraph{Confounding.} A confounder $Z$ causes both $X$ and $Y$, fabricating an $X$--$Y$ association even when $X$ does not cause $Y$. Household \textbf{wealth} is the canonical example here: wealthier districts install more toilets ($\uparrow$sanitation) \emph{and} feed children better ($\downarrow$stunting), so sanitation and stunting correlate without a direct causal link. Table~\ref{tab:corr} confirms the mechanism quantitatively.

\paragraph{The causal ladder.} We climb from association to a defensible effect:
\begin{enumerate}[leftmargin=1.4em,itemsep=1pt]
\item \textbf{Structure learning} --- the PC algorithm (constraint-based conditional-independence tests) infers a DAG (Fig.~\ref{fig:dag}), identifying which variables to adjust for.
\item \textbf{OLS with state fixed effects} --- $Y=\beta X+\gamma^\top Z+\sum_s \delta_s \mathbb{1}[\text{state}=s]+\varepsilon$; read $\hat\beta$ holding wealth and state constant.
\item \textbf{Double Machine Learning} --- partial out confounders from both treatment and outcome with ML and regress the residuals (Frisch--Waugh--Lovell) with cross-fitting, robust to nonlinear confounding.
\item \textbf{Propensity-score matching} --- match high- vs.\ low-$X$ districts on $\hat e(Z)=\Pr(X{=}1\mid Z)$ and compare outcomes.
\item \textbf{Sensitivity (E-value)} --- the minimum strength an unmeasured confounder would need with both $X$ and $Y$ to nullify the effect; large $=$ robust.
\end{enumerate}

\begin{figure}[H]\centering
\figimg{causal_dag}{0.62\textwidth}
\caption{PC-learned causal DAG over NFHS indicators (treatments, mediators, confounders, outcomes).}
\label{fig:dag}
\end{figure}

\paragraph{Results.} Standardized effects across estimators are summarized in Table~\ref{tab:eff} and Figs.~\ref{fig:eff}--\ref{fig:scatter}. The decisive pattern: effects that \emph{collapse toward zero} under adjustment were confounded; effects that \emph{survive} are likely causal.

\begin{table}[H]\centering\small
\begin{tabular}{lcccc l}
\toprule
Relationship & Naive $r$ & OLS+FE & Double-ML & PSM & Verdict\\
\midrule
sanitation $\to$ stunting & $-0.51$ & $-0.15$ & $-0.11$ & $-0.01$ & \textbf{confounded}\\
clean fuel $\to$ stunting & $-0.42$ & $-0.16$ & $-0.14$ & $-0.06$ & mostly confounded\\
schooling $\to$ child marriage & $-0.65$ & $-0.18$ & $-0.17$ & $-0.40$ & \textbf{likely causal}\\
ANC4 $\to$ institutional birth & $+0.60$ & $+0.16$ & $+0.13$ & $+0.36$ & \textbf{likely causal}\\
overweight $\to$ high blood pressure & $+0.57$ & $+0.29$ & $+0.35$ & $+0.44$ & \textbf{causal}\\
\bottomrule
\end{tabular}
\caption{Naive vs.\ adjusted standardized effects. Sanitation$\to$stunting vanishes; schooling, ANC4, and overweight effects persist.}
\label{tab:eff}
\end{table}

\begin{figure}[H]\centering
\figimg{causal_effects}{0.83\textwidth}
\caption{Correlation $\to$ causation across four estimators. Bars collapsing to $0$ were confounded; bars that hold are likely causal.}
\label{fig:eff}
\end{figure}

\begin{figure}[H]\centering
\figimg{causal_scatters}{0.83\textwidth}
\caption{Naive vs.\ wealth-adjusted scatters; the apparent sanitation--stunting slope flattens once wealth is accounted for.}
\label{fig:scatter}
\end{figure}

\paragraph{The result that connects to facilities.} The same machinery shows facility \emph{count} is only weakly linked to outcomes after adjustment. Hence ``build more facilities'' is not automatically the fix; the binding lever is often demand-side (schooling, antenatal care). This is precisely the honesty a ``for good'' planning tool owes its user.

\paragraph{Interactive exploration.} The DAGs are explorable in-app through a force-directed causal-graph studio, with \texttt{d3.js} \emph{inlined} into the Databricks App to satisfy its content-security policy (a CDN-loaded library is blocked). Six modules---the full NFHS map, the wealth-confounding fork, the trust evidence$\to$verdict graph, the care-pathway co-occurrence graph, supply-vs-demand, and the maternal/adolescent chain---offer draggable nodes, zoom, and per-edge evidence on hover.

\section{Statistical ML and geometric deep learning}\label{sec:ml}
Beyond linear adjustment we ran: \textbf{GAM + quantile regression} for nonlinear dose--response and high-burden-district (tail) effects (Fig.~\ref{fig:gam}); \textbf{multilevel/hierarchical} models with empirical-Bayes shrinkage to stabilize small-sample districts; penalized regression with \textbf{stability selection}; and \textbf{geometric deep learning}---a spatial GNN with positional encoding and a Tobler's-Law smoothing penalty---benchmarked against baselines (Fig.~\ref{fig:gnn}). Moran's $I$ is used as a spatial-autocorrelation \emph{diagnostic}, not as proof of a spatial process.

\begin{figure}[H]\centering
\begin{minipage}{0.49\textwidth}\centering
\figimg{gam_quantile}{\textwidth}
\caption{GAM + quantile regression.}\label{fig:gam}
\end{minipage}\hfill
\begin{minipage}{0.49\textwidth}\centering
\figimg{gnn_vs_baseline}{\textwidth}
\caption{Spatial GNN vs.\ baselines.}\label{fig:gnn}
\end{minipage}
\end{figure}

\section{Geospatial methods and referral}
We geocode \emph{any} city/district to a centroid from the $165{,}627$-office India-Post directory (coordinates bounded to India, $6\text{--}38^\circ$N $\times\,67\text{--}98^\circ$E, to reject corrupt values) and rank facilities evidence-first, then by great-circle distance
\[
d=6371\cdot\arccos\!\big(\sin\phi_1\sin\phi_2+\cos\phi_1\cos\phi_2\cos(\lambda_2-\lambda_1)\big)\ \text{km}.
\]

\paragraph{Specialty co-occurrence (``you may also need'').} Beyond the requested specialty $A$, we rank related specialties $B$ by the empirical conditional probability that a facility offering $A$ also offers $B$,
\[
P(B\mid A)=\frac{\#\{\text{facilities offering both }A\text{ and }B\}}{\#\{\text{facilities offering }A\}},
\]
materialized in \texttt{gold\_specialty\_cooccur} over all facilities (kept where $\#A\ge 20$, $\#\{A,B\}\ge 10$ and $P(B\mid A)\ge 0.15$; $6{,}214$ ordered pairs). For instance $P(\text{Internal Medicine}\mid\text{Cardiology})=0.93$ and $P(\text{Orthopedic Surgery}\mid\text{Cardiology})=0.73$. This is a supply-side \emph{bundling} signal, surfaced as a soft suggestion, not medical advice.

\paragraph{Care pathways: cause $\to$ related needs.} A patient routed to a specialty \emph{because of} a condition often needs care for \emph{co-occurring} conditions. We ground this in district-level correlations among the seven patient-facing NFHS indicators (\texttt{gold\_condition\_corr}, all pairs, $n\approx 705$). For a chosen concern we report sibling conditions within the same clinical cluster --- child nutrition (stunting, wasting, child anaemia), metabolic (overweight, high blood sugar), maternal (ANC4, institutional birth) --- annotated with the measured Pearson $r$. Child stunting travels with child anaemia ($r=+0.33$) and wasting ($r=+0.25$); a stunting referral to a paediatrician/child-nutritionist therefore also flags anaemia screening (paediatrics $+$ pathology) and nutrition care nearby, each labelled a \emph{maybe} with an explicit \textbf{population correlation $\neq$ individual diagnosis} caveat. This is where the causal layer (\S\ref{sec:causal}) meets the referral layer: the correlations are real and shown, the confounders (e.g.\ wealth) are examined separately, and the clinical action is \emph{offered, never asserted}.

Roadmap: 2SFCA accessibility scoring to weight reachable supply within a travel radius.

\section{Honesty, uncertainty, and limitations}
Every claim cites the facility's own text; confidence is always shown; the assistant refuses to invent numbers, answering only from grounded gold-table data. Limitations: NFHS-5 is \emph{cross-sectional} (one wave), so our adjusted associations are the strongest single-wave causal claim---two waves would enable difference-in-differences. Scraped supply data carries visibility bias. Suppressed NFHS values are MNAR and are flagged, never imputed.

\section{Conclusion}
facilitiesHelp.io turns $10{,}088$ messy records into decisions a non-technical planner can trust, across all four challenge tracks plus two bonus tracks. The technical core is honesty: grade on cited evidence, communicate uncertainty, separate a real care gap from missing data, and distinguish a correlation from a cause. The hardest part of ``for good'' data work is not prediction---it is telling the user what the data can and cannot support.

\appendix
\section{Data dictionary (gold tables)}
\small
\begin{tabular}{ll}
\toprule
Table & Purpose\\
\midrule
\texttt{silver\_facility} & deduped facility; evidence text; PIN$\to$district\\
\texttt{gold\_facility\_capability\_trust} & 22 acute capabilities + contradiction check ($75{,}335$)\\
\texttt{gold\_facility\_specialty} & every listed specialty graded ($2{,}580$; $117{,}993$)\\
\texttt{gold\_facility\_contact} & address, phone, web, beds, doctors, equipment\\
\texttt{gold\_all\_gaps},\,\texttt{gold\_*\_district\_gaps} & care-gap classification ($2{,}518$ caps)\\
\texttt{gold\_referral} & distance-rankable referral index\\
\texttt{gold\_specialty\_cooccur} & specialty co-occurrence $P(B\mid A)$ ($6{,}214$ pairs)\\
\texttt{gold\_condition\_corr} & NFHS condition correlations for care pathways\\
\texttt{gold\_nfhs},\,\texttt{gold\_pin\_state},\,\texttt{gold\_specialty\_counts} & analysis aggregates\\
\texttt{app\_user\_actions} & persisted overrides/shortlists/scenarios/chats\\
\bottomrule
\end{tabular}

\vspace{6pt}
\noindent\textbf{Live app:} \url{https://virtue-desk-7474648370355072.aws.databricksapps.com}\\
\textbf{Code:} \url{https://github.com/ckwan48/databricks-hackathon-2026}\\
\textbf{Stack:} Databricks Free Edition --- Unity Catalog, serverless SQL Warehouse, Delta medallion, Databricks Apps (Streamlit), Model Serving (Llama 4 Maverick), AI/BI Genie, Databricks agent skills.

\end{document}
